% !TEX root = ../main.tex

\chapter*{前言}
\addcontentsline{toc}{chapter}{前言}

相应结集便于记诵，却不宜当故事书顺读。
研究本必须尊重经号与全覆盖；通读本若仍沿用原短题并列，
则「如来第一」之后接「离贪法第一」，义虽各有所据，气脉已断。

\vspace{0.8em}
\noindent 本册换一问法：\textbf{罗什若为汉地读者重头译出阿含义理，会怎样编排？}
答曰：先安正信，次观五阴、六入、缘起，再明食受谛见与道品，后示修证，兼取善说对机——
篇题自拟，重复熔裁，使如一卷论书可从首读至尾。

\vspace{0.8em}
\noindent 章次：开经·正信 $\rightarrow$ 观五蕴 $\rightarrow$ 观六入 $\rightarrow$ 观缘起 $\rightarrow$
食受谛见 $\rightarrow$ 修道 $\rightarrow$ 修证 $\rightarrow$ 善说与对机。

\vfill
\begin{flushright}
《杂阿含经·仿罗什风新译》项目组\\
2026
\end{flushright}

\clearpage
